% Generated draft: requires mathematical and source review.
\documentclass[11pt]{article}
\usepackage[margin=25mm]{geometry}
\usepackage{fontspec}
\setmainfont{Latin Modern Roman}
\usepackage{amsmath,amssymb,mathtools,unicode-math}
\setmathfont{Latin Modern Math}
\usepackage{enumitem,needspace,xurl,hyperref}
\setlength{\parindent}{0pt}
\setlength{\parskip}{5pt}
\setlength{\emergencystretch}{3em}
\newcommand{\N}{\mathbb{N}}
\newcommand{\Z}{\mathbb{Z}}
\newcommand{\Q}{\mathbb{Q}}
\newcommand{\R}{\mathbb{R}}
\newcommand{\C}{\mathbb{C}}
\newcommand{\F}{\mathbb{F}}
\newcommand{\A}{\mathbb{A}}
\newcommand{\PP}{\mathbb{P}}
\newcommand{\E}{\mathbb{E}}
\newcommand{\eps}{\varepsilon}
\newcommand{\dd}{\,\mathrm{d}}
\begin{document}
\section*{UnsolvedMath ID: 800297}
\textbf{Source problem code:} catalogue entry
\subsection{Definitions and mathematical statement}
Let \(M\) be an interval or a circle, and let \(\operatorname{Homeo}(M)\) denote its group of homeomorphisms. In real projective coordinates, a projective transformation has the form
\[
x\longmapsto \frac{ax+b}{cx+d},\qquad a,b,c,d\in\R,\quad ad-bc\ne0,
\]
interpreted on its domain in the real projective line. A piecewise-projective homeomorphism agrees with projective transformations on its pieces. The supplied record does not specify the permitted partitions, the projective transformation subgroup, or the interval and endpoint conventions; these must be fixed before the class of admissible actions is fully determined.

For a specified smooth structure on \(M\), write \(\operatorname{Diff}^{1}(M)\) for the group of bijections whose maps and inverses are continuously differentiable, with the applicable boundary convention if \(M\) has endpoints. An action \(G\leq\operatorname{Homeo}(M)\) is topologically conjugate to a group of \(C^{1}\) diffeomorphisms if
\[
\exists h\in\operatorname{Homeo}(M)\quad
\forall g\in G,\qquad h\circ g\circ h^{-1}\in\operatorname{Diff}^{1}(M).
\]
The same homeomorphism \(h\) must conjugate every element of \(G\).

\textbf{Classification question (Navas, 2017).} Determine which groups \(G\leq\operatorname{Homeo}(M)\) consisting of piecewise-projective homeomorphisms satisfy this simultaneous conjugacy condition, for interval and circle actions. The supplied statement imposes no finite-generation hypothesis or explicit orientation restriction. This formulation concerns the given action; merely admitting some other abstractly isomorphic \(C^{1}\) action does not express the stated conjugacy requirement. The unresolved conventions above are part of the formulation requiring review.
\subsection{Short English statement}
Classify the piecewise-projective actions on an interval or circle that can be changed into \(C^{1}\) diffeomorphism actions by one common homeomorphic change of coordinates. The source leaves the interval, endpoint, partition, and projective-transformation conventions unspecified.
\subsection{Sources}
\begin{itemize}
\item Supplied local UnsolvedMath registry record: permanent ID \(800297\), exact source problem code \texttt{catalogue entry}, title \emph{Group Actions on 1-Manifolds --- Question 12}; collection: AMR Open Problem Lists. The record attributes the problem to Navas (2017) and labels it \texttt{open}; this is a supplied status label, not an independently verified mathematical or literature assessment.
\item Navas, \emph{Group Actions on 1-Manifolds: A List of Very Concrete Open Questions} (2017), source item recorded as \texttt{Question 12 (\#11)}. Supplied source URL: \url{https://arxiv.org/abs/1712.06462}.
\item The supplied provenance records \texttt{Accessed: 2026-10-16} and extraction method \texttt{arxiv\_latex}. This access date and provenance are supplied curation information; no independent source access or verification is asserted here.
\end{itemize}
\end{document}
